\section{Synthesis and Future Research Plan}\label{sec:plan}

\subsection{Claim--evidence synthesis}

The three completed strands answer different parts of one programme.
\Cref{tab:claim-evidence} separates observations from claims that are already
quantitatively supported and from work that remains open.

\begin{table}[H]
\centering
\caption{Current claim--evidence status.}
\label{tab:claim-evidence}
\small
\begin{tabularx}{\linewidth}{L{0.23\linewidth}L{0.29\linewidth}Y}
\toprule
Claim & Present evidence & Status and next requirement\\
\midrule
CNNs contain a representation-transition regime
& Dataset-wide PCA-RGB depth patterns; random control; semantic-maturity
trajectory
& Plausible and partly quantified; requires frozen, basis-invariant transition
selector across seeds and datasets.\\
\midrule
Late robust feature geometry localizes objects
& ResNet paired tests; frozen adaptive-$K$ validation on four held-out CNN
families after ResNet/ConvNeXt development
& Quantitatively supported on current ImageNet-derived cohorts; requires
cross-dataset and expanded faithfulness confirmation.\\
\midrule
The J1 geometry package benefits student learning
& J1 exceeds H9 across three seeds; implementation-forensics audit; external
single-seed benchmarks
& Consistent with a benefit in the current ConvNeXt ImageNet-100 setting;
component causality remains open, and J1 is practically tied with ReviewKD.\\
\midrule
The same phenomenon extends to transformers
& No closed evidence chain
& Open investigation across several architecture and representation
candidates.\\
\bottomrule
\end{tabularx}
\end{table}

This framing is intentionally conservative. The PhD contribution does not
depend on every method winning every benchmark. It depends on establishing
which geometric properties are reproducible, what question each map answers,
and when those properties produce a useful intervention.

\subsection{Remaining CNN work}

\textbf{The immediate priority is to complete the \methodname{} paper and the
\jmethod{} paper.} This requires closing the current evidence rather than
continuing unbounded formula search. For the PCA-RGB study, this means freezing the
display protocol and transition selector, quantifying wrong-prediction
alignment over the complete error set, and separating class count from other
dataset differences. For \methodname{}, it means a prospective cross-dataset
study with the mode-count selector frozen before segmentation evaluation,
expanded faithfulness and sanity tests, and a fixed semantic atlas containing
failures. For J1, it means replicating the strongest external baselines,
quantifying the J1--ReviewKD difference, reporting resource cost, and testing
generality beyond one teacher--student setting.

\subsection{Transformer investigation across a broad design space}

After the CNN closure, a transformer literature review will begin with a broad
comparison of the
objects explained by attention rollout, gradient-weighted attention,
relevance propagation, token-transformation methods, and feature-formation
methods \cite{abnar2020attentionflow,chefer2021transformer,
xie2023vitcx,wu2024tokentm,mehri2025libragrad,bousselham2025legrad}. The review
will cover flat and hierarchical architectures, supervised and self-supervised
representations, and several candidate explanatory objects. It will then define
a small reproducible baseline set and identify what can be compared fairly
with the CNN representation study.

The first pilot will not assume that a CNN feature grid has a direct
transformer equivalent. Candidate objects include patch tokens before and
after attention, residual-stream states, value vectors, attention-output
features, and hierarchical stage outputs. Candidate architectures may include
a flat supervised Vision Transformer (ViT), a data-efficient variant, and a hierarchical transformer,
but the final selection will depend on reproducible probes and computational
feasibility.

The pilot has three gates:

\begin{enumerate}[leftmargin=*]
    \item \textbf{Measurement gate:} identify a basis-invariant signal of
    token representation change that is stable across seeds.
    \item \textbf{Semantic gate:} show that the candidate transition relates
    to class or object structure rather than attention visualization alone.
    \item \textbf{Utility gate:} demonstrate either improved explanation
    faithfulness or transferable geometry before expanding to larger foundation
    models.
\end{enumerate}

These gates prioritize the most promising evidence path after the broad survey;
they do not presuppose a single architecture or explanatory object. If PCA-RGB
is unstable for token states, the project can use centred kernel alignment
(CKA), linear probes,
affinity change, token clustering, or class-prototype geometry. If
\methodname{} does not transfer directly, that negative result will identify
which CNN assumptions---local spatial grids, convolutional common modes, or
Euclidean channel geometry---do not survive token mixing.

\subsection{Two-year schedule}

\begin{figure}[H]
\centering
\resizebox{\linewidth}{!}{\begin{tikzpicture}[x=1.35cm,y=0.58cm,font=\scriptsize]
\foreach \x/\lab in {0/{Sep--Dec 26},1/{Jan--Apr 27},2/{May--Aug 27},
                       3/{Sep--Dec 27},4/{Jan--Apr 28},5/{May--Aug 28}} {
  \draw[gray!40] (\x,0) -- (\x,-7);
  \node[rotate=45,anchor=west] at (\x+0.08,0.15) {\lab};
}
\draw[gray!40] (6,0) -- (6,-7);
\foreach \y/\lab in {
  1/{CNN evidence and confirmation},
  2/{Transformer literature and pilots},
  3/{Transformer transition study},
  4/{Explanation-method development},
  5/{Geometry-transfer generalization},
  6/{Thesis integration and writing}} {
  \node[anchor=east] at (-0.15,-\y) {\lab};
  \draw[gray!18] (0,-\y-0.35) -- (6,-\y-0.35);
}
\fill[stageblue] (0,-1.25) rectangle (1,-0.75);
\fill[stageblue] (1,-2.25) rectangle (2.15,-1.75);
\fill[stagegreen] (2,-3.25) rectangle (3.15,-2.75);
\fill[stagegreen] (3,-4.25) rectangle (4.15,-3.75);
\fill[stageorange] (4,-5.25) rectangle (5.15,-4.75);
\fill[gray!45] (4.75,-6.25) rectangle (6,-5.75);
\node[diamond,draw,fill=placeholderorange,minimum size=5pt] at (1,-1) {};
\node[diamond,draw,fill=placeholderorange,minimum size=5pt] at (4.05,-4) {};
\node[diamond,draw,fill=placeholderorange,minimum size=5pt] at (5.9,-6) {};
\end{tikzpicture}
}
\caption{Provisional two-year research plan. Diamonds mark the confirmation,
transformer-method, and thesis-integration milestones. The literature-and-pilot
phase surveys several candidates before evidence selects the most productive
direction for deeper study.}
\label{fig:gantt}
\end{figure}

The schedule in \cref{fig:gantt} is organized into six four-month blocks:

\begin{description}[leftmargin=2.6cm,style=nextline]
    \item[Sep--Dec 2026] Complete the confirmation report and finish the
    \methodname{} and \jmethod{} manuscripts, including their remaining
    confirmatory experiments.
    \item[Jan--Apr 2027] Conduct the transformer explainability literature
    review and reproduce a small baseline set.
    \item[May--Aug 2027] Test representation-transition measures on selected
    token-based architectures.
    \item[Sep--Dec 2027] Develop and evaluate one transformer-compatible
    explanation direction if the measurement and semantic gates pass.
    \item[Jan--Apr 2028] Test geometric transfer or cross-architecture
    generalization, conditioned on the earlier findings.
    \item[May--Aug 2028] Consolidate final experiments, integrate thesis
    chapters, and close evidence gaps.
\end{description}

\subsection{Publication plan}

Three intended publications align with the research questions:

\begin{enumerate}[leftmargin=*]
    \item a \methodname{} paper on adaptive robust feature modes for
    class-agnostic CNN explanation;
    \item a \jmethod{} paper on selective hierarchical, variance-balanced
    geometry distillation, after benchmark replication and generality tests;
    and
    \item a transformer representation-transition paper, with its exact method
    determined by the literature review and pilot gates.
\end{enumerate}

\textbf{The first two papers are the next concrete stage of the PhD.} They are
independent enough to publish separately but are linked in the thesis through
representation geometry. The third should test the
scope of the thesis hypothesis rather than introduce an unrelated transformer
method.

\subsection{Risks and contingencies}

The major scientific risk is that the apparent transition does not admit one
universal layer selector. In that case, the thesis can characterize
architecture-dependent transition regimes rather than force universality.
Another risk is that \methodname{} remains strong for localization but weak
for class causality. The contribution should then be positioned explicitly as
class-agnostic object-relevant representation explanation. If J1 remains tied
with ReviewKD, its contribution must rest on verified selective geometry,
causal ablations, and transfer conditions rather than a state-of-the-art claim.
Finally, if transformer geometry differs fundamentally, the thesis can explain
that boundary and develop architecture-specific representations under one
evaluation framework.

\subsection{Professional development and reporting requirements}

The final confirmation submission will include the required academic-integrity
and GenAI-use statements, research training record, publication record, and
resource plan.

\evidenceplaceholder{PROFDEV-TAB-01: researcher training}{
Insert completed Monash Doctoral Program activities, research-methods or ethics
training, conference participation, teaching-related development where
relevant, and planned training for the next reporting period.}
